\section{Introduction}\label{sec:introduction}

Count data exhibiting both spatial dependence and an excess of zeros are pervasive in the applied sciences.
Emergency-department admissions across municipalities, recorded offences across neighbourhoods, firm start-up counts by county, patent citations at the regional level, and disease incidence in small areas all share a common statistical signature: many spatial units record a zero outcome, the non-zero counts display overdispersion relative to the Poisson, and neighbouring units tend to report similar values.
Despite the empirical ubiquity of this data structure, the econometric literature has addressed its three defining features---spatial spillovers, zero-inflation, and overdispersion---largely in isolation or in pairwise combinations, leaving a methodological gap that this paper aims to fill.

In the spatial econometrics tradition initiated by \citet{anselin1988} and surveyed comprehensively by \citet{lesage_pace2009} and \citet{arbia2006,arbia2024}, the workhorse specification for modelling cross-sectional dependence is the spatial autoregressive (SAR) model $\by = \rho\bW\by + \bX\bbeta + \varepsilon$, where $\bW$ is a row-standardised spatial weights matrix and $\rho$ captures endogenous interaction effects among neighbouring units.
For continuous outcomes, maximum-likelihood and generalised-method-of-moments estimators are well established \citep{lee2004,kelejian_prucha1998,kelejian_prucha1999}.
Extending the SAR logic to count data is substantially more challenging because the conditional mean is linked to the linear predictor through a nonlinear (log or exponential) function, and the endogeneity of $\bW\by$ interacts with the discreteness of the support.
\citet{lambert2010} proposed a two-step limited-information maximum likelihood (LIML) estimator for a spatial lag Poisson model, providing a structurally consistent framework that has since become the main reference for SAR-type count models.
However, the Lambert--Brown--Florax estimator handles neither zero-inflation nor overdispersion, two phenomena that routinely co-occur in spatial count data.

On a separate track, the zero-inflated Poisson (ZIP) and zero-inflated negative binomial (ZINB) models, pioneered by \citet{lambert1992} and studied extensively by \citet{cameron_trivedi2013} and \citet{hilbe2011}, accommodate excess zeros through a two-component mixture: a degenerate point mass at zero (the ``structural'' or ``inflation'' component) and a standard count distribution (the ``count'' or ``sampling'' component).
These models have become standard in biostatistics and health economics but, in their classical form, assume independence across observations.
When applied to spatially referenced data, the residuals typically exhibit positive autocorrelation, leading to efficiency losses and potentially biased inference on covariate effects \citep{arbia2006}.

Attempts to marry zero-inflation with spatial structure have taken two main routes.
The spatial-statistics approach introduces conditionally autoregressive (CAR) or geostatistical random effects into one or both components of the ZIP/ZINB model, typically within a Bayesian framework.
Notable examples include the \texttt{CARBayes} package of \citet{lee2013}, the INLA-based disease-mapping models of \citet{asmarian2019}, and the spatial Poisson hurdle specification of \citet{neelon2013}.
While these methods elegantly capture residual spatial heterogeneity, they model dependence as a latent effect rather than as a structural, endogenous interaction.
Consequently, they cannot identify or quantify the direct and indirect (spillover) effects that are central to the spatial econometrics programme \citep{lesage_pace2009,lesage_pace2014}.
The spatial-filtering approach of \citet{metulini2018} and \citet{griffith2003} offers an intermediate solution by projecting spatial dependence onto a set of Moran eigenvectors included as regressors in a standard ZIP/ZINB; this avoids the endogeneity problem but sacrifices the structural interpretation of $\rho$ as a measure of spatial contagion.

No existing framework unifies all four desiderata---endogenous spatial lag, zero-inflation, overdispersion, and unobserved unit-level heterogeneity---within a single, frequentist, estimable specification.
This paper introduces the \emph{Zero-Inflated Spatial Autoregressive} (ZISAR) model to close this gap.
Our contributions are four.
First, we propose a zero-inflated negative binomial model in which the latent linear predictor of the count process follows a simultaneous autoregressive structure: $\boldsymbol{\eta} = (\mathbf{I} - \rho\bW)^{-1}(\bX\bbeta + \bu)$, where the spatial multiplier $(\mathbf{I} - \rho\bW)^{-1}$ propagates covariate effects across the spatial network and unit-level Gaussian random effects $\bu$ capture residual heterogeneity.
The general specification allows a separate SAR process in the zero-inflation component, estimated through a reduced-rank field on a lattice of knots; a restricted version confines spatial dependence to the count predictor.
The model nests several important special cases: setting $\rho=0$ recovers the standard ZINB with random effects; removing random effects and zero-inflation collapses to the aspatial Poisson.

Second, the SAR-on-$\boldsymbol{\eta}$ formulation avoids the endogeneity of observed spatial lags that plagues SAR-on-$\by$ specifications \citep{lambert2010}, enabling proper maximum-likelihood estimation.
We maximise the Laplace-approximated marginal likelihood via Template Model Builder \citep[\texttt{TMB};][]{kristensen2016}, which exploits automatic differentiation and sparse-matrix algebra to scale to $n = O(10^4)$.
This estimation strategy shares the scalability concerns of \citet{arbia_nardelli2024} while delivering exact (up to the Laplace order) marginal inference.

Third, we establish the consistency and asymptotic normality of the Laplace-approximate MLE under regularity conditions analogous to those in \citet{lee2004} and \citet{newey_mcfadden1994}, and derive closed-form expressions for direct, indirect, and total marginal effects that account for both the zero-inflation mixture and the spatial multiplier, adapting the \citet{lesage_pace2009} framework to the nonlinear ZINB setting.

Fourth, we release the estimator, the effects decomposition and prediction in the open-source \texttt{R} package \texttt{zisar} \citep{zisar_pkg}, and demonstrate the model on a Monte~Carlo study and on an application to the $32{,}322$ small areas of England, which the sparse log-determinant of Section~\ref{sec:est_laplace} makes feasible at full scale.
In the adoption of domestic solar photovoltaic systems during the boom of the Feed-in Tariff both parts of the restricted model are needed---spatial dependence in the counts and structural zeros where the housing stock leaves no roofs to equip---and the effects decomposition attributes the larger part of every total effect to spillovers.

The remainder of the paper is organised as follows.
Section~\ref{sec:literature} surveys the existing literature on spatial count models, zero-inflated spatial specifications, and spatial filtering approaches.
Section~\ref{sec:model} introduces the ZISAR model and states its key properties.
Section~\ref{sec:estimation} details the Laplace-approximate ML procedure and derives the asymptotic results and effects decomposition.
Section~\ref{sec:montecarlo} reports a Monte~Carlo study comparing ZISAR with competing specifications and the power of tests for spatial dependence.
Section~\ref{sec:solar} applies the restricted model to the adoption of domestic solar photovoltaic systems in the small areas of England during the boom of the Feed-in Tariff and reports its direct and indirect effects.
Section~\ref{sec:conclusion} concludes and discusses extensions, including the incorporation of temporal dependence.
Proofs, the simulation checks of the general model and implementation details are collected in the Appendix.
